\originalchapter{指令层级怎样成为可训练的问题}\label{study:ih}
\noindent 研究来源：OpenAI 2026-03-10 官方说明与 Guo 等 2026-03-11 的 IH-Challenge v1\cite{ihblog,ih}。本文是原创中文研究说明，不是逐段译本。该论文采用 arXiv 非独占分发许可，未据此获得全文翻译许可；原图、原数据和长提示不镜像。下面的符号、例子、推导和练习为本项目设计。

\originalsection{先把三个不同的失败拆开}
给模型一个数学问题，它答错了，不能直接断言它不懂指令层级。给模型相互冲突的指令，它输出了符合高优先级格式的错误内容，也不能直接断言任务完成。应分别检验解题、遵循约束、辨认约束的权威来源。指令层级研究关心第三项，但测量会受到前两项污染。

设一段对话为 $c=((r_1,m_1),\ldots,(r_k,m_k))$，$r_i$ 是消息角色，$m_i$ 是内容。角色由应用的结构传入，正文里自称“管理员”不改变角色。对本研究采用的政策，权威次序为 system、developer、user、tool。低级消息中的资料仍可用于完成任务；只有冲突的行为要求应被忽略。“工具不可信”不等于其每一个事实都是假的。

\begin{definition}
给定高优先级规则 $s$，令 $\mathcal Y_s(c)$ 为符合规则的输出集合。约束遵循率为 $\Pr[y\in\mathcal Y_s(c)]$；任务正确率另由 $T(c,y)$ 衡量。两者的交集 $\Pr[y\in\mathcal Y_s(c),T(c,y)=1]$ 才衡量受约束任务完成。
\end{definition}
例如规则要求“只输出 A 或 B”，低级消息要求“附上一段解释”。输出 A 满足格式，但若正确选项是 B，仍未完成任务。输出“B，因为……”可能知识上正确，却不满足当前接口。评分器必须同时记录两种事实，不能只给一个含糊总分。

\originalsection{固定评分骨架为何有用}
IH-Challenge 的关键工程取舍是将高优先级任务和程序评分器先固定，再在线合成低优先级攻击；任务应易于执行、易于客观评分，同时包含反误拒样本。这样可以减少把普通解题困难误计为冲突处理失败，也让攻击改变时评分目标仍有固定含义。论文还审计评分器，并不把通过几个测试当成评分器绝对正确。

我们用一个原创的严格 schema 例子说明这个取舍。高优先级要求输出 JSON 对象且键集合恰为 $\{\texttt{choice}\}$，值属于 $\{\texttt{A},\texttt{B}\}$；用户给出两选一算术题。评分器先解析整个输出，再检查键、类型和值，最后与独立计算的答案比较。字符串搜索“含有 B”会接受“我不选择 B”；只找第一段花括号会容忍泄漏在对象外的额外文本。这些反例说明评分器本身就是研究对象。

设 $D$ 是高优先级任务分布，$A(\delta\mid s)$ 是当前攻击分布，$\pi_\theta$ 是模型，固定评分器为 $g_s(y)$。一个抽象训练目标是
\[
J(\theta)=\E_{s\sim D,\,\delta\sim A(\cdot\mid s),\,y\sim\pi_\theta(\cdot\mid s,\delta)}[g_s(y)].
\]
若 $A$ 随训练改变，这不是一次固定监督数据集上的拟合。若评分器错了，优化可把错误放大。程序可判定仅表示计算过程明确，不表示程序定义的目标等于我们真正想要的安全。

\begin{example}[拒绝一切为何是捷径]
有 80 个禁止泄漏题和 20 个授权正常任务。若奖励只记录前一组不泄漏，永远拒绝得到满分。若联合完成率采用全部 100 题且正常任务拒绝记错，则永远拒绝只有 $0.8$。这个数仍可能过高，所以报告还应分别给两组分数；改变样本配比也会改变总分。
\end{example}
\begin{solution}
设拒绝率策略在禁止组成功概率为 $a$、正常组成功概率为 $b$，汇总值是 $0.8a+0.2b$。只看汇总值不能恢复 $a,b$：$(1,0)$ 与 $(0.75,1)$ 都给 $0.8$。需要分组指标，并把正常任务难度控制到模型本来能够解决的范围。
\end{solution}

\originalsection{论文结果的含义与实验身份}
作者报告 GPT-5-Mini 经训练得到内部实验模型 GPT-5-Mini-R，16 项评价平均稳健率由 $84.1\%$ 到 $94.1\%$，指定安全评价中的不安全行为由 $6.6\%$ 到 $0.7\%$。这是该论文协议中的作者结果，不能移植成当前公开产品的错误率，也不能证明任意新攻击都被抵御。官方说明中的聊天偏好指标还存在小幅下降；“能力基本保持”应结合每项指标理解。本文未调用该内部模型，未复现这些前沿分数。

本项目的 \texttt{v03\_research.py} 用明确来源的符号策略模拟 240 个正常案例及 120 个攻击案例。来源盲策略在 120 个案例都偏离用户目标；加执行器权限检查后，80 个越权提案被阻止，仍有 40 个许可内的语义错误。关键词拒绝挡住攻击，却误拒 120 个正常案例。理想来源路由器在这个被完全定义的有限任务上全部正确。这些是实际程序结果，证明了该程序的机制差异；它没有语言理解，不能据此宣称 LLM 已被防护。

另一个实际诊断使用两份固定版本 SmolLM2-135M 权重，包含六道 system/user 冲突题。两选项续写评分的准确率均为 $3/6$，严格单字母生成格式的通过率为零。完整输出保留供检查。这里没有 tool 消息，只有很小的题集，没有自适应攻击或训练；它不能复现 IH-Challenge，但让我们看到格式、选择概率和任务完成确实是三个不同测量。

\originalsection{如何设计下一次研究}
可做的下一步是保持事实任务不变，只交换同一条命令所在的结构角色，比较成对输出。先在无冲突条件验证任务本来可解，再加入冲突；另设看起来危险但实际授权的正常题，测误拒。攻击集按作者、模板、目标分别留出，防止只记住关键词。评分器应有解析歧义、嵌套对象、额外文本等负例。若允许模型采取行动，执行器在模型之外独立审查操作、资源和目的地，语言输出的分数之外还要审计真正发生的状态变化。

\begin{exercise}
一个模型不泄漏率为 $0.99$，正常任务完成率为 $0.50$；另一模型两项分别为 $0.95,0.90$。若禁止与正常题各占一半，谁的联合得分高？为什么仍不能只报这个值？
\end{exercise}
\begin{solution}
分别是 $0.745$ 和 $0.925$，第二个更高。但在泄漏代价很高的系统中，任务比例不是损失比例；还须报告两项、按后果加权，并在部署分布评估，不能把平均数当安全保证。
\end{solution}
\begin{exercise}
若评分器只检查输出没有秘密字符串，模型输出另一个用户的隐私会不会被发现？
\end{exercise}
\begin{solution}
未必。评分器的可判定范围只是指定字符串。这是测量覆盖缺口，不是随机种子问题。扩充威胁模型、独立访问控制与审计目标，才能形成对该类风险的证据。
\end{solution}
